\chapter{Continuum correlations and local fields}
\label{ch13:continuum}

The ultraviolet limit leaves a theory on each fixed physical torus.
We now let its period grow. The construction may select a subsequence
of volumes, but it does not select a new subsequence of ultraviolet
cutoffs for the correlations. Local fields are then defined by their
responses in all separated correlations. Polynomial labels that give
the same responses represent the same field.

We use the full fixed-volume limit of
Theorem~\ref{ch12:uv-theorem}, the volume-uniform finite-difference
bounds of Theorem~\ref{ch11:all-period-differences}, and the equation comparisons and the comparisons using separating families in Chapter~\ref{ch11:insertions}.
Compactness gives the volume extraction; the separate equation
inclusion and rank bounds identify its local fields.

\section{One subsequence of volumes}
\label{ch13:extraction}

Recall that $\Conf_n$ is the open set of pairwise distinct points in
$(\RR^4)^n$, without any restriction on their coordinate projections.
Let $p(k)$ be any strictly increasing sequence of volume indices.
On every compact subset of $\Conf_n$, the periodic kernels
$M^{\rm UV}_{p(k),v}$ are eventually defined with separation bounded
below independently of $k$. Each of their derivatives is uniformly
bounded there. The cutoff threshold used to prove the bound may
depend on the fixed volume; after taking that ultraviolet limit it
disappears from the bound.

Arzel\`a--Ascoli and diagonal extraction
\cite[Theorem~7.25]{Rudin1976} give one refinement $P(k)=p(q(k))$
and a multilinear family $S_n$ such that
\begin{equation}
 M^{\rm UV}_{P(k),v}\longrightarrow S_n(v;\cdot)
                    \quad\hbox{in }C^\infty_{\rm loc}(\Conf_n).
 \label{ch13:volume-limit}
\end{equation}
The diagonal family is countable: arities, degree lists, basis-label
lists, derivative lists, and compact exhaustions of the configuration
spaces. Multilinearity extends its conclusions to all labels.
Integration by parts identifies the extracted derivative limits with
the derivatives of $S_n$. No diagonalization over all smooth test functions is
needed. For each fixed separated test function $f$ the order of limits is
\begin{equation}
 \lim_{k\to\infty}\lim_{N\to\infty}
 a_N^{4n}\sum_y f(a_Ny)M_{N,P(k)}(v;y)
                   =\int f(x)S_n(v;x)\,dx.
 \label{ch13:ordered-limit}
\end{equation}
The inner limit is the whole cutoff sequence. Every further volume
refinement preserves \eqref{ch13:volume-limit} and
\eqref{ch13:ordered-limit}.

Permutation symmetry, conjugation, the distinguished unit and
hypercubic covariance pass from the finite correlations by these
limits. Translation by an arbitrary real vector is obtained by
approximating it on the same lattice; the uniform first-difference
bound controls the rounding error. Full orthogonal covariance will
require a separate argument.

\section{Separated cells and errors in local observables}
\label{ch13:cells}

The chessboard seminorm of Section~\ref{ch01:reflection} bounds an
observable by its repeated reflected copies in a tiling. For local
observables $F_1,\ldots,F_n$ interior to distinct cells of the same
translated tiling $\mathcal T$, the chessboard inequality gives
\begin{equation}
 \left|\left\langle\prod_iF_i\right\rangle_{N,p}\right|
                   \le\prod_i\|F_i\|_{\mathcal T}.
 \label{ch13:chess}
\end{equation}
It is the lattice inequality already cited from
\cite{FrohlichIsraelLiebSimon1978}, with the constant observable in every unused cell.
The estimate does not grow with the distances between the cells.

Here is the required geometric placement. For a tiling of physical
side $a$, the translations in one coordinate bringing one of $n$
anchors within distance $w$ of a boundary have total length at most
$2nw$ modulo $a$. If $2nw<a$, a translation avoids all these strips;
a lattice translation does so for sufficiently fine cutoff. Choose
one in each coordinate. Anchors with pairwise maximum-coordinate
distance greater than $a$ cannot share a cell. Hence if the centered
local bounds apply at scale $k$ with margin $3a_k$, and
$6na_k<a_m$, anchors separated by more than $a_m$ admit distinct
side-$a_m$ cells with the required margins.

We shall use an error family $E_{N,p}$ with the following property:
for any fixed arity and sufficiently small fixed cell side, its
chessboard seminorm is at most $\eta_{N,p}\to0$ in the ordered
limits, uniformly over placements satisfying a specified interior
margin. The cell side can be chosen after the arity and a requested
lower scale, before the cutoff. The spectator insertions have
uniform centered seminorm bounds on sufficiently fine fixed scales.
Place the error and all spectators in distinct cells, choosing a
still finer scale for the spectators' margins. Equation
\eqref{ch13:chess} yields
\begin{equation}
 \left|\left\langle E_{N,p}(y_i)
                  \prod_{h\ne i}O_{N,p}(v_h,y_h)\right\rangle\right|
               \le C(v,\delta)\eta_{N,p},
 \label{ch13:vanishing-error}
\end{equation}
for configurations separated by $\delta>0$. The weighted number of
lattice points in the fixed compact support of a test function is bounded. The same
error therefore vanishes after separated smearing. This conclusion
uses the stated order of choices for the error norm; pointwise
smallness at one configuration would not suffice.

\section{Fields as equivalence classes of labels}
\label{ch13:quotient}

A label includes its degree space: the same formal polynomial in
$V_d$ and $V_e$ can have different cutoff normalizations. For
$c\in V_d$, a fixed spectator list $v$, and a separated smooth test function
$f$, define the one-slot response
\[
 L_{v,f}(c)=\int f(x_0,x)S_{n+1}(c(x_0),v(x))\,dx_0\,dx.
\]
Let
\begin{equation}
 N_d(S)=\bigcap_{v,f}\ker L_{v,f},\qquad F_d=V_d/N_d(S).
 \label{ch13:null}
\end{equation}
The intersection permits every finite spectator list, of arbitrary
degrees. Smoothness makes this equivalent to vanishing in every
pointwise separated correlation. Permutation symmetry ensures that
nullity has the same meaning in every slot. The quotient $F_d$ is
finite dimensional, but its dimension has not yet been identified.

We use the bounded equivariant idempotents $\Pi_{d,N,p}$ of
Section~\ref{ch11:uniform-differences}, with kernel
$\mathsf Z_{N,p,d}P_{N,p,d}S_{t_N}R_d\mathcal I_d$.
Choose a common ordered limit $\Pi_d$ after a coefficient refinement. More
precisely, diagonal extraction over the countably many degrees and
volumes gives one increasing sequence $N=n(j)$ along which the bounded
coefficient arrays converge at each fixed volume. A further volume
refinement makes those limits converge as well. All coefficient
limits below use these refined ordered tails. This auxiliary selection
does not change the already existing full cutoff limits of the
correlations. Write
$r_d=\dim\mathcal M_d^{\rm Max}$ for the dimension of the free Maxwell
quotient in Chapter~\ref{ch03:maxwell}. The construction gives:
\begin{enumerate}
\item the ranks of $\Pi_{d,N,p}$ are at most $r_d$;
\item every bounded family in $\ker\Pi_{d,N,p}$ has vanishing
separated responses in the ordered limits;
\item there are finite matrices of limiting separated responses,
evaluated on comparison labels, that approach an invertible
$r_d\times r_d$ free matrix of correlation pairings.
\end{enumerate}
The first statement is exact rank equality, since every map defining
the finite equation space is invertible. For the second,
Lemma~\ref{ch11:uniform-sweep} makes its chessboard seminorms tend
to zero; bounded rows have uniformly subpower inverse terminal
coordinates, and \eqref{ch13:vanishing-error} supplies every spectator
response. For the third, Proposition~\ref{ch11:finite-detector-input}
gives the actual finite-cutoff matrices with error tending to zero
as their comparison depth increases. Include their finitely many
bounded label columns at each of the countably many comparison depths
in the same coefficient refinement. At each fixed depth, their
ordered limits preserve the matrix estimate by multilinearity and
the uniform correlation bounds. An injective rectangular matrix of correlation pairings
has an invertible square minor; retain those rows.

\begin{proposition}[Identification by two rank bounds]
\label{ch13:rank}
For this selected continuum family,
$N_d(S)=\ker\Pi_d$ and $\dim F_d=r_d$.
\end{proposition}
\begin{proof}
For fixed $c$, the bounded family $(1-\Pi_{d,N,p})c$ has vanishing
responses. Taking limits gives $L_{v,f}((1-\Pi_d)c)=0$ for all
$v,f$, so $\ker\Pi_d\subset N_d(S)$. Converging idempotents have
eventually constant ranks: if their norm difference is less than
one, each is injective on the other's range. Thus
$\dim F_d\le\operatorname{rank}\Pi_d\le r_d$.
Each comparison matrix factors through $F_d$, since its rows are
separated responses. Matrices sufficiently close to the invertible
free matrix of correlation pairings have rank $r_d$. Hence $\dim F_d\ge r_d$.
Equality of dimensions and the kernel inclusion prove the result.
\end{proof}

\subsection{Representatives chosen before the spectators}

Approximate inclusion and derivative labels need not stay bounded.
The following finite-dimensional observation explains the normalized
null estimate required to replace them by fixed vectors. It also
avoids choosing a separate representative for every test function.

Let $V$ be finite dimensional, $\Pi_u\to\Pi$ idempotents, and
$L_{i,u}\to L_i$ linear functionals, along a proper filter of indices.
Assume $\ker\Pi=\bigcap_i\ker L_i$. If vectors $x_u$ satisfy
\begin{equation}
 L_{i,u}(x_u)\to t_i,\qquad
 L_{i,u}\left(\frac{x_u-\Pi_ux_u}{1+\|\Pi_ux_u\|}\right)\to0
                              \quad\hbox{for every }i,
 \label{ch13:normalized-null}
\end{equation}
there is one $z\in\operatorname{ran}\Pi$ with $L_i(z)=t_i$ for
all $i$. To see this, put $b_u=(1+\|\Pi_ux_u\|)^{-1}$ and
$z_u=b_u\Pi_ux_u$. A common cluster point $(b,z)$ satisfies
$b+\|z\|=1$, $\Pi z=z$ and $L_i(z)=bt_i$. If $b=0$, then
$z\in\ker\Pi\cap\operatorname{ran}\Pi$, a contradiction.
The vector $z/b$ is the desired representative. Compactness is used
only in one fixed finite-dimensional space; all response identities
hold at the same cluster point because they hold along the filter.

In the application, the filter is generated by ordered tails:
$p=P(k)$ with $k\ge k_0$, and $N=n(j)$ with $j\ge j_0(k)$.
The normalized null condition is available for every admissible
coefficient family in the sense of Section~\ref{ch11:uniform-differences}.
Indeed, if $x_{N,p}$ has subpower inverse terminal coordinates, then
$x/(1+\|\Pi x\|)$ has the same property, whereas
$\Pi x/(1+\|\Pi x\|)$ is bounded and therefore admissible.
Their difference lies in the exact finite equation kernel.
Lemma~\ref{ch11:uniform-sweep}, followed by
\eqref{ch13:vanishing-error}, proves
\eqref{ch13:normalized-null} for all spectators on the same tail.
This argument applies even when the original $x_{N,p}$ are unbounded.

\section{Inclusions and derivatives}
\label{ch13:calculus}

For each basis label $c\in V_d$, the protected trial construction
of Section~\ref{ch11:trial-descent} gives approximate representatives
in $V_e$ for $e\ge d$ and in $V_{d+1}$ for a physical lattice
difference. To check the inclusion trial, multiply the original
cutoff observable by $a_N^e$; its local representative is
$a_N^{e-d}$ times the degree-$d$ section, so its derivative and
value bounds remain subpower. For a difference, the extra factor
$a_N$ cancels the inverse step, and the fixed translated-stencil
cost is bounded by the coupling-threshold requirements used in
Theorem~\ref{ch11:all-period-differences}.
Trial descent yields admissible coefficients and vanishing chessboard seminorm errors, uniformly in period. The preceding paragraph gives
their normalized null estimates.
All representatives and comparison errors are chosen before
the spectators and test functions.
Applying the preceding
finite-dimensional observation gives fixed linear maps on labels,
$\widehat\iota_{d,e}$ and $\widehat\partial_{d,\mu}$, such that
\begin{align}
 S(\widehat\iota_{d,e}c(x_0),v(x))&=S(c(x_0),v(x)),\label{ch13:inclusion-response}\\
 S(\widehat\partial_{d,\mu}c(x_0),v(x))
                      &=\partial_{x_0^\mu}S(c(x_0),v(x)).
 \label{ch13:derivative-response}
\end{align}
For the derivative target, exact torus summation by parts is
\[
 a_N^{4n}\sum_y f(a_Ny)\Delta_{N,i,\mu}M_{N,p}(y)
 =a_N^{4n}\sum_y
 \frac{f(a_Ny-a_Ne_{i,\mu})-f(a_Ny)}{a_N}M_{N,p}(y).
\]
The quotient on the right tends uniformly to $-\partial_{i,\mu}f$;
the zero-order separated bound controls its error. This identifies
the distributional derivative target. Smoothness gives
\eqref{ch13:derivative-response} pointwise.

These maps descend to unique maps
$\iota_{d,e}:F_d\to F_e$ and $\partial_\mu:F_d\to F_{d+1}$
with the displayed responses. They preserve null classes by their
definition. The inclusions are injective and satisfy
\[
 \iota_{d,d}=1,\qquad
 \iota_{e,f}\iota_{d,e}=\iota_{d,f},\qquad
 \partial_\mu\partial_\nu=\partial_\nu\partial_\mu.
\]
Each identity follows by testing against every spectator and using
\eqref{ch13:null}; the derivative maps also commute with inclusions.
Conjugation and hypercubic covariance descend in the same way.
The union of the $F_d$ is therefore a filtered space of local fields,
with derivatives raising the filtration by one. No multiplication of
two fields at the same point has been defined by this construction.

\section{Reflection positivity and separated power bounds}
\label{ch13:positivity}

Let a word mean a finite list of labels and positions, with correlation
denoted $S(w)$. Write $w^*$ for conjugation of its labels and reversal
of its order. Time is coordinate four. Let $\theta w$ reflect that
coordinate and act on its labels by the corresponding signed
transformation, as well as applying $*$. For words with strictly
positive times, lattice bond reflection positivity passes to
\[
 \sum_{i,j}\overline{c_i}c_j S(\theta w_i\,w_j)\ge0.
\]
It also passes for joint smooth test functions supported on positive-time
separated configurations. Indeed their lattice supports stay away
from both reflection planes on the ordered tail. The lattice plane
$x^4=a_N/2$ tends to zero, and the uniform separated bound controls
the moving test functions. Taking the two limits preserves positivity of each
finite matrix. Pointwise positivity follows by concentrating smooth
bumps. This uses the finite lattice result of
Section~\ref{ch01:reflection}; it asserts no positivity for unspecified
extensions at coincident points.

For later use take the centered chessboard estimate
$\|O\|_k\le C a_k^{-(d+1)}$, uniform on a bounded set of labels
of degrees at most $d$ and every sufficiently fine fixed level $k$.
It follows by choosing $\epsilon=1$ in
Proposition~\ref{ch13:sharp-centered}, whose proof below uses only
the bound for reflected tilings of Chapter~\ref{ch11:insertions} and the
selected schedule. Slow growth of the block
ratios gives a constant $c>0$ with
\begin{equation}
                         a_{k+1}\ge c a_k^2.
 \label{ch13:scale-thickness}
\end{equation}
The first scale below a length $r\le1$ is consequently at least a
constant times $r^2$. Choose a cell side below $r/2$, then a centered
scale below that side divided by $6(n+1)$. The latter is at least
$c_n r^4$. The separated-cell estimate yields
\begin{equation}
 |S_n(v;x)|\le C_{n,v}\,r^{-4(d+1)n},
 \qquad \min_{i\ne j}|x_i-x_j|_\infty\ge r.
 \label{ch13:power}
\end{equation}
The constant is independent of the diameter of the configuration.
Applying the same argument to the fixed degree-raised labels in
\eqref{ch13:derivative-response} gives a polynomial separated bound
for every derivative.

\section{Repeated reflected words}
\label{ch13:repeated}

The preceding bound cannot be used with arity tending to infinity
without controlling its constants. We need a more specific estimate.
Let $w$ have $n\ge1$ insertions, degrees at most $d$ and coefficient
norms at most $R$,
pairwise separation at least $r$, and times in $[r,h-r]$, where
$0<r\le1$. We call $r$ its buffer in the slab of width $h$.
For a positive-time exterior word $v$, set
\[
 A_wv=w\,\tau_{se_4}v,\qquad w^\dagger=\tau_{se_4}\theta w,
 \qquad \|v\|_{\rm OS}^2=S(\theta v\,v).
\]
At this point $A_w$ acts only on words; boundedness on a Hilbert
space is not assumed.

\begin{proposition}[Uniform bound for repetitions]
\label{ch13:word-bound}
Under the centered chessboard bounds and their interior-cell
comparison used above, choose $0<s\le h$ so that $w$ still has
buffer $r/2$. There are $C=C(d,R)$ and $\kappa=4(d+1)$ such that,
for every exterior word $v$,
\[
 \|(A_{w^\dagger}A_w)^jv\|_{\rm OS}^2
       \le B_v\,[C^n(n/r)^{\kappa n}]^{4j},\qquad j\ge0.
\]
Here $s,C,\kappa$ are chosen before $v$, and $B_v<\infty$ is
independent of $j$.
\end{proposition}
\begin{proof}
Choose a cell scale $a_m<r/32$ as the first permitted such scale,
and $s=a_m\lfloor h/a_m\rfloor$. Equation
\eqref{ch13:scale-thickness} gives $a_m\ge c_1r^2$.
Choose a centered scale $a_k<\alpha(r/n)^2$, with $20\alpha<c_1$.
Again taking the first such scale gives
$20na_k<a_m$ and $a_k\ge c_2(r/n)^4$.

The expanded reflection form contains $4jn$ main insertions and two
exterior copies. Since $s$ is a multiple of $a_m$, the main points
have at most $2n$ residue patterns modulo the side-$a_m$ tiling,
independently of $j$. The strip count therefore places all main
points in distinct cells with margins $5a_k$, using constants
independent of the number of repetitions.

Choose finer aligned cell and centered scales for the finitely many
exterior points, with exterior cell side at most $a_k$.
Their two copies have times $\pm(2js+t)$, so a
choice depending on their minimum positive time and internal
separation works for every $j$. A further translation smaller than
the exterior cell side adjusts their margins while leaving the main
margins at least $3a_k$. The odd integer ratios preserve the fine
cell grid under coarse reflections. Apply the chessboard inequality
first to the coarse cells, regarding all exterior factors in one
coarse cell as a single observable, then to the fine cells for those
factors. If $X$ bounds the exterior centered seminorms, the result is
\[
 X^{2|v|}(C_0a_k^{-(d+1)})^{4jn}.
\]
Take $B_v=\max(1,X^{2|v|})$ and use the lower bound on $a_k$.
Every use of the lattice inequality is at fixed $j$ before taking
the ordered limits. A cutoff threshold uniform in $j$ is unnecessary.
\end{proof}

This estimate is the input for bounded slab insertions in the
reflection Hilbert space. Its derivation precedes that operator
construction and does not use orthogonal covariance or a phase
decomposition. The equation comparison, admissible trial encoding,
estimates for finite separating families and centered estimates all refer to the same selected
family. Full orthogonal covariance and products at coincident points
remain separate constructions.

% !TEX root = ../monograph.tex
% Fragment for Chapter13: insert after the coarse separated bound.
\section{Sharp short-distance bounds for a fixed number of insertions}
\label{ch13:sharp-section}

For a fixed number of insertions the loss in
\eqref{ch13:power} can be reduced to an arbitrarily small excess over
the sum of their degrees. This uses the normalized integral over the reflected tiling,
not the stronger Gaussian comparison used for the separating families. We keep
all incoming joint components with insertions; replacing them by a product of
terms of source degree one would invalidate the estimate.

Write $\tau_m=\log(1/a_m)$, so
$\widehat\beta_m=\bar\beta+\kappa_G\tau_m$. Let $C_j\ge1$ include
the inverse-column, finite-section and source constants at stage $j$.
The stage ratio and the lower bound on the reference coupling may be chosen to satisfy
\begin{equation}
 \frac{\log C_j}{\log \mathscr L_j}\longrightarrow0,\qquad
 B_{i+1}^2\le2B_i^2,\qquad
 \frac{(L_iB_i^2)^4 q_{i+1}^{\rm site}}{\tau_{i+1}}
                          \longrightarrow0.
 \label{ch13:sharp-schedule}
\end{equation}
Here $B_i$ is the support-block radius, and $q_i^{\rm site}\ge1$ is
the per-site normalization exponent of \eqref{ch11:replicated-bound};
it is not a density with insertions or the coefficient of the coupling constant renormalization.
For the path-action representation of the density one can use
$q_i^{\rm site}\le C_{J(i)}(1+\log\widehat\beta_i)$.
Indeed the exponent without its constant term is at most $C_{J(i)}V$. Set
$\rho=\widehat\beta_i^{-1/2}$ and restrict every link to
$\|U_e-1\|_G\le\rho/4$, using the matrix norm \eqref{ch01:matrix-norm}.
On this event the exponent is at least $-C_{J(i)}V$ and all large-field
indicators vanish. Exponential summability in the path length bounds the
leading positive path quadratic by $C_{J(i)}\widehat\beta_i\rho^2V$.
Signed regular costs and the absolute
activity site sum satisfy the bounds \eqref{ch04:activity-bounds}. The lower bound for the partition function
of Section~\ref{ch11:replicated} therefore adds only
$C_{J(i)}V+2D_GV\log\widehat\beta_i$ to the exponent, where $D_G$ is an ambient
real matrix dimension for the gauge group.

These requirements are compatible with sufficiently large lower bounds on the reference coupling. At each fixed
stage $B_i=O_j((\log(2+\tau_i))^2)$ and
$q_i^{\rm site}=O_j(1+\log(2+\tau_i))$, so the last quotient in
\eqref{ch13:sharp-schedule} tends to zero at its fixed ratio. Choose
the ratio after $C_j$, then increase the lower bound on the reference coupling until that quotient is at
most $1/j$ for both adjacent-stage choices. The same choice gives
the displayed bound on consecutive $B_i$,
$\log L_i=o(\tau_i)$, and $\log C_{J(i+1)}=o(\tau_i)$
at transitions. Choose each ratio against both adjacent-stage column
constants, which are fixed beforehand. These are conditions on the single selected
schedule; no new schedule is chosen after fixing a correlation.

\subsection{A normalized local insertion}

Fix a scaling-dimension bound $d$ and a ball of radius $R$ in the
coefficient space of $V_d$. Set
\[
 E_k=C_{d,R}^{\rm sec} C_{J(k)}\prod_{i<k}C_{J(i+1)},\qquad
 \omega_k=a_k^d\widehat\beta_k^{-w_d}E_k^{-1}.
\]
The fixed section constant $C_{d,R}^{\rm sec}$ includes the coefficient radius and all its signed and
conjugated images. Choose the fixed integer $w_d$ large enough to absorb
the finite powers of the inverse small-field threshold in the degree-$d$
section and its counterpart with a prescribed large-field set, as well as the source attachment bound. Multiplication of the single-insertion inverse bounds gives the
factor $E_k$; the scaling normalization is exactly $a_k^{-d}$.
Thus multiplication of the \emph{full} insertion by $\omega_k$
makes the exponential bound for its regular term of source degree one
less than three and the bound for its large-field term of source degree one
less than one. Its remaining normalized size can be placed below the
attachment bound by increasing $w_d$.

The bounds in Chapter~\ref{ch11:insertions} for terms of source degree at least two retain every
incoming source term and every original eligibility event. Decreasing all source
weights only decreases their sums of absolute coefficients. The source attachment
inequalities at the intervening scales propagate this chosen small
weight; they require the same inverse-column and threshold powers already
included in $E_k,w_d$. Consequently \eqref{ch11:root-bound}, with
$x=3,y=1$, gives
\begin{equation}
 \|O_{N,p}^{(d)}(c)\|_{\mathrm{cell}(m)}
 \le \omega_k^{-1}
           \exp\bigl((a_m/a_k)^4q_k^{\rm site}+6\bigr).
 \label{ch13:sharp-root}
\end{equation}
It holds for all sufficiently large $N$ uniformly in period and in the coefficient
ball, whenever the reflected insertion points at the stopping scale have spacing
at least $B_k^2$. Their spacing stays adequate at finer scales because
$L_i\ge3$ whereas $B_{i+1}^2/B_i^2\le2$.

The logarithm of $E_k$ is $o(\tau_k)$: split its sum into a finite
initial segment and a tail where
$\log C_{J(i+1)}\le\epsilon\log L_i$, then let $\epsilon\downarrow0$.
The separate terminal column is absorbed by the same ratio bounds and
$\log L_{k-1}=o(\tau_k)$. Also
$\log\widehat\beta_k=o(\tau_k)$. Therefore
\begin{equation}
                        6-\log\omega_k=d\tau_k+o(\tau_k).
 \label{ch13:weight-size}
\end{equation}
This is the reason the prefactor from the reflected tiling does not introduce a fixed
extra scaling degree.

\subsection{The stopping scale and centered estimate}

Choose $k=k(m)$ as the first scale $k\ge m$ with
$B_k^2\le a_m/a_k$. It exists because $B_k$ grows subexponentially
in $\tau_k$, whereas $a_m/a_k=e^{\tau_k-\tau_m}$. By minimality,
\[
 (a_m/a_k)^4q_k^{\rm site}
       \le (L_{k-1}B_{k-1}^2)^4q_k^{\rm site}=o(\tau_k).
\]
If the first scale already qualifies the same conclusion follows from
the fixed-stage bound. Since $q_k^{\rm site}\ge1$,
\[
 0\le\tau_k-\tau_m=\log(a_m/a_k)
       \le(a_m/a_k)^4q_k^{\rm site}=o(\tau_k).
\]
It follows that $\tau_k/\tau_m\to1$.
Combining this with \eqref{ch13:sharp-root} and
\eqref{ch13:weight-size} proves the following statement.

\begin{proposition}[Sharp centered bound]
\label{ch13:sharp-centered}
For every scaling-dimension bound $d$, coefficient radius $R$ and $\epsilon>0$, there
is $C_{d,R,\epsilon}$, fixed before the cell scale $m$, such that
\[
 \|O_{N,p}^{(d)}(c,y)\|_{\mathrm{cell}(m,y)}
        \le C_{d,R,\epsilon}a_m^{-d-\epsilon},\qquad\|c\|\le R.
\]
For each fixed $m$ this holds for all sufficiently large cutoffs,
uniformly in $p,y$ and in signed and conjugated images. The insertion
is interior in that cell on the same tail.
\end{proposition}
\begin{proof}
For all sufficiently large $m$, the logarithm of the right side of
\eqref{ch13:sharp-root} is at most $(d+\epsilon)\tau_m$ by the
preceding limits. For the finite remaining set of $m$, use its chosen
fixed stopping scale and take the maximum of the finite constants.
The fixed link support becomes interior at every fixed cell
scale in continuum units as $N\to\infty$. Translation covariance makes the lower bound on $N$
independent of $y$.
\end{proof}

\subsection{Arbitrary separated points}

The slow ratio condition gives a sharper estimate of the first scale below a length
than \eqref{ch13:scale-thickness}. For any fixed $c>0$, let $m(r)$
be the first scale with $a_m<cr$. Then
\begin{equation}
            \frac{\log(1/a_{m(r)})}{\log(1/r)}\longrightarrow1
                         \quad(r\downarrow0).
 \label{ch13:first-below}
\end{equation}
Indeed $a_{m-1}\ge cr>a_m=a_{m-1}/L_{m-1}$, and the logarithm of
the overshoot is at most $\log L_{m-1}=o(\tau_{m-1})$.
Any fixed finite iteration of such choices of a first scale has the same property.

For a fixed number $n$ of insertions and mutually separated insertion points, first choose a
cell side below a fixed small multiple of their minimum separation $r$.
Translate its boundaries away from all insertion points. Then choose a centered
scale smaller than this side divided by $6(n+1)$.
The interior-cell comparison and the chessboard inequality apply.
Both chosen scales have logarithm $(1+o(1))\log(1/r)$ by
\eqref{ch13:first-below}. Apply Proposition~\ref{ch13:sharp-centered}
with a sufficiently small excess to each degree separately. We obtain,
for every $\epsilon>0$,
\begin{equation}
 |S_n(A_1,\ldots,A_n;x)|
    \le C_{\boldsymbol A,n,\epsilon}
             r^{-\sum_i d_i-\epsilon},\qquad 0<r\le1,
 \quad A_i\in F_{d_i},\quad\min_{i\ne j}|x_i-x_j|_\infty\ge r.
 \label{ch13:sharp-separated}
\end{equation}
The same bound is uniform over bounded sets of fixed raw representatives.
Its constant is independent of the diameter of the configuration.
For fixed derivatives, replace the differentiated field by its degree-raised
label; their total order is added to $\sum_i d_i$.


% !TEX root = ../monograph.tex
\section{Bounded real approximations for separated words}
\label{ch13:bounded-approximations}

The later phase argument needs real approximants bounded uniformly in
the cutoff. Compactness of the link space at each cutoff is insufficient:
the physical normalization of a local field grows as the mesh shrinks.
We instead stop at one fixed coarse level and pull back a bounded coarse
function through the probability block kernels.

We record the precise marked estimate used in this construction. It is
the local-trial comparison for reflected tilings, with every term involving at least two insertions
retained, not only its one-column coefficient estimate. At a coarse
level $h$, let $J_h$ map the coefficients of a trial in the chosen basis
to the coefficients of its selected level-$h$ polynomial. Write $C_h$
for the prescribed coefficient-column bound,
$Q_h$ for the ordinary norm bound of a normalized local trial, and
$E_h$ for its complete bound, including the complementary large-field
part. For the selected trajectory these bounds and all fixed coefficient
products are subpower in $a_h^{-1}$. The last-step endpoint estimate
and chessboard seminorm estimate for the residual required here are
\begin{align}
 \|J_h-I\|_{1\to1}&\le
           2(1+C_h)Q_hL_{h-1}^{-1/8}\le\tfrac14,
 \label{ch13:coarse-matrix}\\
 \|\text{trial residual}\|_{\mathcal C_k}
       &\le C_{k,m,d}E_h e^{-(\tau_h-\tau_m)/8}.
 \label{ch13:coarse-residual}
\end{align}
Here $k,m$ are fixed before $h$; $\mathcal C_k$ is the centered
cell tiling at level $k$ from Section~\ref{ch13:cells}. Both estimates must concern the history restricted from
the actual cutoff $N$, uniformly in $N\ge h$ and the physical period.

The first inequality is the endpoint comparison used in
Proposition~\ref{ch11:equation-comparison}. Its final $1/4$ is an
additional finite stage requirement: require $L_{h-1}^{1/8}\ge8(1+C_h)Q_h$. For the normalized small-field polynomial the local smooth cutoff
is identically one on the ordinary chart, so its ordinary derivative
constant is fixed before this block choice. Its global value is controlled by
the square-root large-field bound, by a larger coupling threshold as in
Section~\ref{ch11:equation-calibration}. This does not replace $E_h$ by
$Q_h$ in the second estimate.

The second inequality requires more than a small coefficient linear in the sources.
We prove its complete marked version next. The source weight is increased
across successive scales; it is not increased by its entire final factor
on an already formed coarse component.

\begin{lemma}[The coarse-trial residual on reflected tilings]
\label{ch13:coarse-residual-supplier}
On the selected histories, after the finite coupling-threshold requirements
specified below, \eqref{ch13:coarse-residual} holds with all incoming
regular and large-field terms involving at least two insertions retained. Its constant is independent
of the outer cutoff, the period, and the number of replica cells.
\end{lemma}
\begin{proof}
Fix the degree ceiling before the schedule. At level $i$, choose a
positive number $\zeta_i$ so that a term linear in the sources with complete normalized
size at most $\zeta_i$ fits every local attachment and adjacent
comparison bound of Chapter~\ref{ch11:insertions}. Include the finite
signed symmetry orbit in these bounds. More explicitly, if $D_i$ is
the complete adjacent comparison constant and $u_i>0$ the terminal
attachment bound, require
\[
 C_i\zeta_i\le1,\qquad 2D_i\zeta_{i+1}\le1,\qquad
 e^{81\sigma_i+162}Z_i^{\rm val}\zeta_i/u_i\le1.
\]
Here $Z_i^{\rm val}$ bounds the global section and its complementary
cutoff part; this symbol denotes a bound, not a partition function.
Also decrease $\zeta_i$ to put the incidence sum for terms linear in the sources below the
small constant in \eqref{ch11:higher-induction}. These choices concern
normalized coefficients, before the physical factor $a_i^{-d}$.

These positive upper bounds can be chosen with
$|\log\zeta_i|=O_j(1+\log\widehat\beta_i)$ on every fixed stage.
Indeed the section derivatives, inverse small-field thresholds, and
adjacent comparison constants have finite powers of
$\widehat\beta_i$, with stage-dependent constants. The terminal attachment bound may be taken as $u_i=\tfrac14\widehat\beta_i^{-1/2}$; the original large-field suppression
is reserved separately in the attachment estimate and need not be
inverted to define $u_i$. The factor $e^{81\sigma_i+162}$ is constant
on each stage. Taking the minimum of this finite list of reciprocal
bounds gives such a choice. At transitions between adjacent stages use the minimum
of both lists. All requirements below therefore follow by increasing
the coupling thresholds, since $B_i$ grows as a positive constant times
$(\log\widehat\beta_i)^2$ on a fixed stage.

Choose $0<\phi_i\le1$ such that
\begin{equation}
 \phi_i L_i^{d+1}\max(1,C_{i+1})\zeta_i\le\zeta_{i+1},
 \qquad \phi_i^{-1}\le e^{B_i/128}.
 \label{ch13:residual-height}
\end{equation}
For example take the minimum of one and the ratio in the first
inequality. The logarithm of its inverse is bounded by a finite
stage constant times $1+\log\widehat\beta_i$, plus
$(d+1)\log L_i$, and hence is eventually below $B_i/256$.
Any finite local incidence bound $\delta_i>0$ can simultaneously
be included so that $\delta_i\phi_i e^{B_i/128}\ge1$:
its inverse has at most another fixed power of
$\widehat\beta_i$. This verifies the comparison and attachment bounds as well as the displayed, weaker upper bound on
$\phi_i^{-1}$. These are countably many finite degree requests to
the single schedule of coupling thresholds.

Put $G(h,i)=e^{(\tau_h-\tau_i)/8}$ and choose
\[
 w=\frac{\zeta_m a_m^dG(h,m)}{2E_h},\qquad
 \omega_i^{\rm res}=w\prod_{j=m}^{i-1}\phi_j.
\]
The residual coefficients from \eqref{ch11:descent-error}, with the
fixed harmless column factors included in $E_h$, satisfy
$c_i+|x_i|\le E_h/G(h,i)$. Since the ratio of
$a_i^{-d}/G(h,i)$ at successive levels is $L_i^{d+1/8}$,
\eqref{ch13:residual-height} gives
\begin{equation}
 \omega_i^{\rm res} a_i^{-d}\bigl(c_i+|x_i|\bigr)
       \le \omega_i^{\rm res} a_i^{-d}E_h/G(h,i)\le\zeta_i/2
                  \quad(m\le i\le h).
 \label{ch13:residual-height-bound}
\end{equation}
In particular the negative normalized local insertion used in the subtraction has
normalized cost at most $\zeta_h$ at $h$. Above $h$, its inverse
coefficient bounds grow by at most $C_{i+1}$ per step; the same bound
recurrence controls these factors and the physical factor $L_i^d$.
Thus its terminal source satisfies the joint attachment bound at $N$,
and its complete push to $h$ is available at the weights $\omega_i^{\rm res}$.

Here is the support calculation for its terms involving at least two insertions, including those
formed above $h$. Open the centered cells at $m$ with spacing at least
$B_m^2$. This spacing persists because $L_i\ge3$ and the adjacent
support radii satisfy the schedule of \eqref{ch13:sharp-schedule}.
Apply reweighting on the incoming fine support $X$. If it contains
$r\ge2$ distinct marked ancestors, then
$B_{i+1}r\le64c_{B_{i+1}}(X)$ by \eqref{ch11:separated-cover}.
Since $B_i\le B_{i+1}$ and $\omega_i^{\rm res}/\omega_{i+1}^{\rm res}=\phi_i^{-1}$, one factor from the incoming exponential support weight gives
\[
 (\omega_i^{\rm res}/\omega_{i+1}^{\rm res})^r e^{-c_{B_{i+1}}(X)}
 \le e^{B_i r/128-c_{B_{i+1}}(X)}
 \le e^{-c_{B_{i+1}}(X)/2}.
\]
To make this argument iterable, let $A_i>4$ be the full exponent
maintained in the norm for terms involving at least two insertions at level $i$, and impose
\[
 10\cdot15^4(A_i+9)\le L_i(A_{i+1}-4).
\]
Reserve four units of the incoming $A_{i+1}$ norm for this weight
increase, site conversion, and attachments. The support contraction
of Chapter~\ref{ch11:insertions} converts the remaining exponent
$A_{i+1}-4$ into the full outgoing exponent $A_i$, including its local
reserves. Thus the estimate does not lose two exponent units at each
successive step. Choosing the bounds for terms linear in the sources small enough makes
\eqref{ch11:higher-induction} preserve the bound one in these full norms
for both arrays of terms involving at least two insertions. This calculation applies to every cardinality and
period. A large-field array is treated as one incoming factor with
its complete original prescribed set of large-field sites, not as a product of its coefficients linear in the sources. The terminal joint construction reserves the original
large-field decay for the complementary factors. Hence there is no
uncontrolled reweighting by $w^r$ at the intermediate level $h$.

At that level insert the positive local trial in every cell.
If $\mathscr A_h$ is the complete pushed negative source, the new polynomial is
\[
 \mathscr A_h\exp_{\rm nil}\Bigl(\sum_v z_vF_v\Bigr).
\]
In its local representation this adds $F_v$ only to the regular
entry linear in the sources. All regular terms involving at least two insertions and all original large-field
terms remain unchanged. The scalar part is exactly the coefficient of the identity in the global exponential.
The nonconstant part multiplied by the smooth cutoff has zero complementary profile because
its small-field restriction is included in its definition.
Linearity makes the new term linear in the sources precisely the residual whose
coefficients satisfy \eqref{ch13:residual-height-bound}; its adjacent
regular and large-field comparisons have bounds at most one by the
choice of $\zeta_i$. Descend this whole polynomial to $m$. Its regular terms linear in the sources have bound at most two, and the
large-field terms linear in the sources have bound at most one, and its two incidence bounds for terms involving at least two insertions are at most one.

Let $\mathsf P_{h,N}$ push densities from level $N$ to level $h$ by
successive applications of \eqref{ch04:block-operator} along the chosen
history. Its adjoint $\mathcal P_{N,h}$ pulls observables back from $h$ to $N$.
Let $\mathcal I$ index the cells of the reflected tiling and let $-U_v$
denote the terminal insertion used for cell $v$. In the product above
the coefficient of $\prod_{v\in\mathcal I}z_v$ at level $h$ is
\[
 \sum_{A\subseteq\mathcal I}
 \mathsf P_{h,N}\left(\rho_N\prod_{v\notin A}(-U_v)\right)
                       \prod_{v\in A}F_v.
\]
Duality and the disjoint-cone identity
\eqref{ch13:cone-factorization} identify its integral with the original
partition function times the expectation of
$\prod_v(\mathcal P_{N,h}F_v-U_v)$.
No product of separately pushed densities for single insertions has replaced it.
Equation~\eqref{ch11:root-bound} now yields
\[
 \|\mathcal P_{N,h}F-U\|_{\mathcal C_k}
 \le w^{-1}\exp\bigl((a_k/a_m)^4q_m^{\rm site}+5\bigr)
 =\frac{2e^{(a_k/a_m)^4q_m^{\rm site}+5}}{\zeta_ma_m^d}
       E_h e^{-(\tau_h-\tau_m)/8}.
\]
The prefactor depends only on the fixed stopping data. This proves
\eqref{ch13:coarse-residual}, once the cones have the room verified
below. All identities and estimates use the original selected history.
\end{proof}

\begin{proposition}[Coarse bounded approximation]
\label{ch13:bounded-singleton}
Use the selected histories with the stage requirements above. Fix a scaling-dimension bound $d$, a real reference coefficient $c$, a
sufficiently fine centered tiling $k$, and $\delta>0$. There are real
observables $\mathcal B_{N,p}(c,y)$ with a finite bound independent of
$N,p,y$, such that
\[
 \|\mathcal B_{N,p}(c,y)-O_{N,p}(d,c;y)\|_{\mathcal C_k(y)}\le\delta
\]
on the ordered cutoff and volume tails. Both observables are interior
to the same cell. The construction includes their signed and conjugated
images with the same bound.
\end{proposition}
\begin{proof}
Fix $m\ge k$ with $B_m^2\le a_k/a_m$. At a finer level $h$, use the
coarse trial
\[
 V_h(v;U)=a_h^{-d}\left(v_0+
       \mathbf1_{\text{small-field stencil}}
                    P_d^{\rm loc}[S_{1/t_h}v^\circ](U)\right).
\]
The scalar part is retained exactly. The full inspected stencil includes
the links needed by the relaxed small-field condition. For fixed $h$ this is
a bounded measurable function, with bound uniform in the outer cutoff
and period: its coefficient range and allowed couplings are compact,
and its section is a fixed polynomial in compact group variables.

Let $\mathcal K_{N,h}(U,dV)$ be the composition of the one-step probability kernels
$k^{(L_i)}_\beta(U,V)\,dV$ of \eqref{ch04:kernel}, in order from
$i=N-1$ down to $h$ and with the coupling of each step, and set
$\mathcal P_{N,h}F(U)=\int F(V)\mathcal K_{N,h}(U,dV)$.
It preserves real values and contracts the supremum norm. Locality here
must include every random variable used to choose a gauge-covariant block average or gauge
frame. Conditional on the preceding link field, use independent primitive
random variables for those choices and for each coarse-link sample.
The full ancestor cone of a stencil contains its fine-link reads and
all these random variables at every intermediate level. For disjoint
cones, conditioning and integrating the independent variables gives
\begin{equation}
 \prod_C\mathcal P_{N,h}F_C
                   =\mathcal P_{N,h}\left(\prod_CF_C\right).
 \label{ch13:cone-factorization}
\end{equation}
This does not assert that coarse links remain conditionally independent
after a shared frame choice has been averaged out. The full cones,
including that choice, must be disjoint. Here is an explicit uniform
radius bound. Let $j_i$ be the smoothing order in the step $i+1\to i$.
In Chapter~\ref{ch04:blocking}, each endpoint frame reads a box of
fine-lattice radius $(j_i+3)L_i$; the other endpoint is at distance
$L_i$, and the connector paths stay in these boxes and the intervening
strip. Thus every input link and every shared frame seed for an output
link lies within physical distance $(j_i+7)a_i$ of its initial endpoint,
in the maximum norm. This bound also includes eligibility conditions and
one boundary-link margin. If the coarse inspected stencil has radius
$r_d$ in level-$h$ units, its full ancestor cone therefore has radius at
most
\[
 r_da_h+\sum_{i=h}^{N-1}(j_i+7)a_i.
\]
The growing smoothing orders require a schedule condition, not merely
finiteness at each step. Impose $(j_i+7)\le a_i^{-1/2}$ at all sufficiently
late levels. This is a permitted finite coupling-threshold requirement: if
$\jmath_q$ is the largest smoothing order used at stage $q$ or its
permitted adjacent pairs, require
$b_q\log3\ge2\log(\jmath_q+7)$ among the coupling-threshold conditions of
Section~\ref{ch10:schedule}. When $J(i+1)=q>1$, that schedule has
$i\ge b_q$ and $a_i\le3^{-i}$, giving the desired inequality.
The finitely many initial levels are irrelevant after increasing $h$.
Since $a_{i+1}\le a_i/3$, the radius is bounded, independently of $N$, by
\[
 r_da_h+\frac{a_h^{1/2}}{1-3^{-1/2}}\longrightarrow0.
\]
Choose $h$ so that this is smaller than the fixed interior margin of
each reflected level-$k$ cell. The cones, including their shared seeds,
then lie strictly inside disjoint cells; sufficiently large physical
periods prevent wrapping identifications.

The bound \eqref{ch13:coarse-matrix} gives
$\|J_h^{-1}\|_1\le2$. Thus the reference insertion with coefficient
$c$ is a linear combination of the trial columns, with sum of absolute
coefficients bounded by
$2\|c\|_1\exp(\sum_{i<h}\log C_{i+1})$. Replace these columns by
$\mathcal P_{N,h}V_h(e_j)$ and call the resulting observable
$\mathcal B'$. Its supremum is finite for fixed $h$ and independent
of $N,p,y$.

Replicate the discrepancies before applying the kernel. At level $h$
the coarse factors multiply the complete source density, including its
terms involving at least two insertions. Equation~\eqref{ch13:cone-factorization} identifies the
resulting residual with the product of the pullback discrepancies.
Lemma~\ref{ch13:coarse-residual-supplier} and recombination give
\[
 \|\mathcal B'-O(c)\|_{\mathcal C_k}
 \le C_{k,m,d,c}\exp\left(\sum_{i<h}\log C_{i+1}\right)
                       E_h e^{-(\tau_h-\tau_m)/8}.
\]
The prefactor is $e^{o(\tau_h)}$, so choose $h$ to make this at most
$\delta$, before choosing the cutoff. For real $c$ take
$\mathcal B=\operatorname{Re}\mathcal B'$; conjugation invariance
and the triangle inequality for the chessboard seminorm preserve the
bound. Translations and the finite symmetry orbit use the same stencils
and constants. Outside the ordered tails set the approximants to zero
if a bound for every index is desired.
\end{proof}

\begin{corollary}[Real word differences]
\label{ch13:bounded-word}
On these selected histories, the bounded-real approximation used for phase
reflection in Section~\ref{ch15:phase-symmetry} holds on a cofinal set of
generic grids. Its supremum bound is independent of copy number and copy
position; approximation thresholds may depend on each fixed pair.
\end{corollary}
\begin{proof}
Choose the grid level after the degree bound and any prescribed lower
level, then fix a generic real word $w$ of length $n$. Include all its
signed and conjugated labels in a finite family. Choose centered cells
with margins fitting between its distinct coordinate projections, and
let $C$ bound the single-insertion seminorms there. Apply
Proposition~\ref{ch13:bounded-singleton} with error $\delta\le1$.
The approximants have seminorm at most $C+\delta$, regardless of their
possibly much larger supremum norms.

For a product of $q$ such separated factors, telescoping and the
chessboard inequality give
\[
 \left|\left\langle\prod_{a=1}^q\mathcal B_a-\prod_{a=1}^qO_a
               \right\rangle\right|
                  \le q\delta(C+\delta)^{q-1}.
\]
Form the difference $g_j$ of the bounded products for a translated word
and its linear time reflection, using the same single-insertion approximation
at every translated copy. If their supremum bounds for single insertions are $M_i$,
then $|g_j|\le2\prod_iM_i=:R_\varepsilon$, independent of all copy
positions and indices. For two sufficiently widely separated copies,
expand their product of differences into four products of $2n$ factors.
Choose $\delta$ with $8n\delta(C+1)^{2n-1}\le\varepsilon$.
The preceding inequality gives the desired pair error. Only finitely
many configurations are involved at each fixed pair, so ordered tails
can be intersected. The empty word difference is zero. Increasing the
initial lower-level request gives the required cofinal grids.
\end{proof}


% !TEX root = ../monograph.tex
\section{Scaling dimension and the physical quotient}
\label{ch18:graded-quotient-section}

An estimate on a polynomial coefficient is useful for local product replacement only if
it retains the degree of the field that coefficient represents.
We derive the required quotient estimate from the triangular marked
transfer. This step uses the corrected physical inclusions and the
normalized-null statement of Section~\ref{ch13:quotient}; those statements
were established using the exact microscopic equation space.

Fix an output scaling dimension $t$. For $v\in V_t$, let
$x_l(v)=(x_{l,\alpha}(v))_\alpha$ be the vector of polynomial coefficients
after transport to level $l$, in the fixed basis of polynomials of
scaling dimension at most $t$. Write $|x|_1=\sum_\alpha|x_\alpha|$. The leading polynomial map at one step has the form
$M_l\mathsf d_l$, where
\[
 (\mathsf d_lx)_\alpha=L_l^{t-|\alpha|}x_\alpha.
\]
Here $|\alpha|$ is the scaling dimension of the basis polynomial, not
its factor count. The leading map fixes the constant and is triangular:
$(M_l)_{\alpha\beta}=0$ off the diagonal whenever
$|\alpha|\le|\beta|$. Both $M_l$ and its inverse have norm at most the
prescribed stage constant $K_l$. The full map differs from it by the
marked remainder bound of Section~\ref{ch11:one-graph}.
Enlarge the prescribed column bounds $C_l$ and block floors so that
\begin{equation}
 64K_l^2\le C_l,C_{l+1},\qquad C_{l+1}\le L_l,
 \qquad 2K_lL_l^{-3/8}(1+C_{l+1})\le1.
 \label{ch18:graded-room}
\end{equation}
These are finite-stage inequalities, with constants fixed before the
block ratio, and the schedule gives
$\sum_{l<m}\log C_{l+1}=o(\tau_m)$.

\begin{proposition}[Graded quotient estimate]
\label{ch18:graded-estimate}
For $0\le i\le t$ there are constants $C_{t,i}(m)$ with
$\log C_{t,i}(m)=o(\tau_m)$ such that the following holds.
Suppose $v_{N,p}\in V_t$ converges on an ordered cutoff and volume
refinement to $v$, and $|x_m(v_{N,p})|_1\le\epsilon$ on that tail.
Then
\begin{equation}
 \operatorname{dist}_{F_t}([v],\iota_{i-1,t}F_{i-1})
       \le C_{t,i}(m)a_m^{-(t-i)}\epsilon,
 \label{ch18:graded-distance}
\end{equation}
where $F_{-1}=0$. Constants are independent of the refinement and the
convergent label family. The norm on $F_t$ is fixed once.
\end{proposition}
\begin{proof}
For $i>0$ let $P$ project onto scaling dimensions below $i$.
Triangularity makes $\ker P$ invariant under $M_l$ and its inverse.
The degree scaling expands low coordinates by at least
$L_l^{t-i+1}$, and high coordinates by at most $L_l^{t-i}$.
The one-step marked bound and \eqref{ch18:graded-room} imply
\begin{align}
 |x_{l+1}|_1\le C_{l+1}|Px_{l+1}|_1
 &\ \Longrightarrow\quad
 |x_l|_1\le C_l|Px_l|_1,\quad
 L_l^{t-i+1}|Px_{l+1}|_1\le C_{l+1}|Px_l|_1,
 \label{ch18:low-cone}\\
 Px_{l+1}=0
 &\ \Longrightarrow\quad
 |x_l|_1\le C_{l+1}L_l^{t-i}|x_{l+1}|_1.
 \label{ch18:high-step}
\end{align}
To check the first implication, put $Y=|P\mathsf d_lx_{l+1}|_1$.
The premise and $C_{l+1}\le L_l$ give
$|\mathsf d_lx_{l+1}|_1\le2Y$ and $|x_{l+1}|_1\le Y$. Applying the inverse leading map
to the low part and absorbing the marked error gives
$Y\le2K_l|Px_l|_1$. The full output is at most $3K_lY$;
\eqref{ch18:graded-room} gives both assertions. When the low part
is zero, the high-degree scaling and $K_l+1\le C_{l+1}$ give
\eqref{ch18:high-step}.

At the microscopic level take the low plane $PV_t$ and transport it.
Equation~\eqref{ch18:low-cone} shows that its projection onto the low
coordinates at every later level is injective, and hence onto that
fixed-dimensional coordinate space. At level $m$ subtract its vector
with the same low coordinates as $x_m(v_{N,p})$. The remaining
coefficient has zero low part and norm at most $(1+C_m)\epsilon$.
Repeat the subtraction after each step towards the fixed reference
level, applying \eqref{ch18:high-step} between subtractions. The sum of
all subtracted vectors still lies in the transported microscopic low
plane. Calling its reference label $y_{N,p}$, we obtain
\[
 |v_{N,p}-y_{N,p}|_1\le
 G_m(1+C_m)(a_{r(t)}/a_m)^{t-i}\epsilon,
 \qquad G_m=\prod_{l<m}(1+C_l)C_{l+1}.
\]
Finite initial transfer and norm-equivalence constants are absorbed in
$G_m$. Its logarithm is $o(\tau_m)$ by the prescribed column bounds.

It remains to identify this microscopic low plane with physical lower
degrees; it is not the naive inclusion of unnormalized polynomials.
The labels $y_{N,p}$ are bounded by the displayed inequality.
The expansion in \eqref{ch18:low-cone}, iterated from the microscopic
level, bounds their microscopic low coefficients by
$a_N^{t-i+1}e^{o(\tau_N)}$. Therefore their degree-$t$ physical
representatives can be written with the normalization $a_N^{-(i-1)}$
and coefficient/derivative growth $e^{o(\tau_N)}$, on the same fixed
stencil. Their global value bounds also have this growth, since the
threshold inverses have subpower growth.

The local-trial comparison of Chapter~\ref{ch11:insertions} matches
these representatives to degree-$(i-1)$ reference labels with vanishing
separated residuals. The matched coefficients need not be bounded.
Apply the normalized-null representative argument of
Section~\ref{ch13:quotient}: division by one plus the norm of their
corrected projection makes the equation remainder uniformly controlled,
and its separated responses vanish. Since $y_{N,p}$ has a bounded
reference limit, that argument gives a physical degree-$(i-1)$ field
with exactly its limiting responses. The all-spectator definition and
corrected inclusion thus put every such limit $[y]$ in
$\iota_{i-1,t}F_{i-1}$. Passing to a convergent subsequence of the bounded
$y_{N,p}$ proves \eqref{ch18:graded-distance}. Extending a tail-defined
bounded family by zero outside its tail changes none of these limits.

For $i=0$ no low plane is removed. The reverse coefficient estimate
directly bounds $|v|_1$ by
$e^{o(\tau_m)}(a_{r(t)}/a_m)^t\epsilon$. Passing to the quotient
proves the same conclusion. All fixed reference spacings are absorbed
in $C_{t,i}(m)$.
\end{proof}

The power $t-i$ in \eqref{ch18:graded-distance} is what recovers a
coefficient of degree $i$ from an auxiliary extraction at a much higher
degree $t$. An ungraded bound on $[v]$ would lose precisely that power.
The proof also explains why identifying the lower field requires all
spectator responses, even though its eventual coefficient is measured
by a fixed finite family of probes.

% !TEX root = ../monograph.tex
% May be placed in Chapter13 after the graded quotient estimate.
% This fragment uses neither phase decomposition nor stress-tensor Ward identities.
\section{Local replacement of a cluster on the lattice}
\label{ch18:lattice-products}

We prove a local replacement estimate before using any phase decomposition.
It will also apply to the undecomposed correlations, before any decomposition into phases, with a fixed exterior word.
Observable source variables and local polynomial extraction are standard
renormalization-group tools; see, for example,
\cite{BauerschmidtBrydgesSlade2015}. The quantitative estimates below concern
the block kernels, activity norms and quotient fields constructed
here; that reference does not give these particular estimates.

For $T\to\infty$, write $f\preccurlyeq e^{cT}$ if, for every
$\varepsilon>0$, eventually $|f|\le e^{(c+\varepsilon)T}$ on an ordered
cutoff and volume tail, uniformly on the specified bounded label sets
and geometries. A factor $e^{o(T)}$ is called subpower in this paragraph.
Choose scales $l=l(T)$ with $\tau_l/T\to\theta$, where
$\tau_l=\log(1/a_l)$. The sequence of scales has $\log L_l=o(\tau_l)$, so
first-crossing scales give any positive $\theta$. For $\theta=0$ one may
first cross $\sqrt T$.

\subsection{Signed local transfer}

A local degree-$t$ density with insertions consists of a column coefficient
$x_l$ and a remainder of norm $r_l$ whose extracted coefficient vanishes. Its dimensionless
normalization obeys
$H_l=L_l^t\mathcal B_l H_{l+1}$; the dimensionful density with insertions has
no additional factor. These conventions are those of
Section~\ref{ch11:one-graph}. Keep the full prescribed large-field set and
its normalization by the square root of the large-field bound in this remainder norm.

The backward comparison of Section~\ref{ch11:trial-descent} applies
on any sufficiently late interval $[k_*,m]$, not only from the original
cutoff scale. In detail, a remainder of size $q$ with vanishing extracted coefficient has
output coefficient and remainder at most $L_l^{-2/5}q$.
The regular term uses its zero radial Taylor polynomial through $t$;
the large-field term retains half the suppression of its original
prescribed large-field set. Subtracting its output column restores the vanishing extracted coefficient and
adds at most $(1+2C_l)L_l^{-2/5}q$. The lower bound on the block ratio makes this at
most $L_l^{-3/8}q$. This recentering is an equality of
densities with insertions, including the terms with prescribed large-field sets, and uses linearity rather
than division by an activity.

Add a free column parameter at the finest scale and choose it so that
the transferred residual coefficient vanishes at $k_*$. Invertibility
of the column map permits this on the same history. The cone calculation
in Section~\ref{ch11:trial-descent} then gives
\begin{equation}
 r_l+|x_l|_1\le4(1+C_m)q_{\rm fin}(a_m/a_l)^{1/8},
 \qquad r_m\le2(1+C_m)q_{\rm fin}.
 \label{ch18:signed-transport}
\end{equation}
Here $q_{\rm fin}$ is the combined original coefficient and remainder
bound, before the affine column subtraction. To recall the numerical
point, with $Y=|\mathsf d_lx_{l+1}|_1$, $U=|x_l|_1$ and
$\epsilon_l=L_l^{-3/8}$ the two inequalities are
$Y\le K_l[U+\epsilon_l(Y+r_{l+1})]$ and
$r_l\le K_lY+\epsilon_lr_{l+1}$.
The cone $32(1+C_l)^2U\le r_l$ gives
$Y\le4K_l\epsilon_lr_{l+1}$ and $L_l^{1/8}r_l\le r_{l+1}$.
It starts with $x_{k_*}=0$. At the finest scale absorb the subtracted
column bound to obtain the second inequality, and iterate for the first.
The sign of the local difference has been retained throughout.

\subsection{Grouping insertion points and retaining their supports}

Use the square-zero source algebra of Chapter~\ref{ch11:insertions}.
The total number of copies may be arbitrary, but a reflected cell
contains a fixed number $n$ of insertion points. A group is a set of
these insertion points with a common representative. Grouping is bookkeeping for
supports and source weights: original variables and the coefficient of every source monomial
remain in the expansion. In particular partial modes of a group are
not discarded, and grouping is not asserted to be an algebra
homomorphism replacing their product by a new variable.

Recall that $c_B(X)$ counts the radius-$B$ balls in a smallest connected covering of $X$, and $\sigma$ is its exponential support weight.
If the support $X$ of a term with insertions has to contain new representatives joined
by at most $h$ covering balls per insertion point, then
\[
 c_B(X\cup\{y_i:i\in I\})\le c_B(X)+h|I|.
\]
Shrinking each affected weight by $\eta e^{-\sigma h}$ absorbs this
support enlargement. The incidence norm contributes a second factor: a new
incidence can be attributed to an old insertion point on its joining path.
If at most $\mathfrak N$ representatives occupy a covering ball, there
are at most $(1+\mathfrak N)n$ assignments per representative.
Take $\eta\le[(1+\mathfrak N)n]^{-1}$. Supports without insertions are
unchanged. This is a termwise estimate on the full
series with insertions, hence is uniform in its number of replica cells.

At each scale retain a maximal $B_l^2$-separated set of representatives,
keeping the representative of the center, and join each discarded
representative to a retained neighbor. A merger decreases the number
of groups; there are at most $n-1$ mergers within a cell.
A deliberately generous joining bound is
\[
 h_l=16[1+L_l(B_l+B_{l+1}+1)^2]^4.
\]
Impose $(\sigma_{l+1}+1)h_l=o(\tau_l)$ by choosing the lower bounds on the reference coupling at each stage.
At fixed stage $h_l$ is a polynomial in $\log\widehat\beta_l$
with fixed ratio, so this is a finite admissible request.
The occupancy factors have the same property. The minimum
$\alpha_j\le1$ of all weight-reduction factors for support enlargement on $m\le l<j$
therefore has $-\log\alpha_j=o(\tau_j)$: split the maximum
cost into a fixed initial segment and a tail bounded by
$\varepsilon\tau_j$, and then let $\varepsilon\downarrow0$.

Between mergers a single group loses at most the factor $6L_l$ in
its weight. Here the incoming regular contribution is carried
\emph{without} Taylor subtraction. Its full factor for the change of cover radius
is $T_l\le L_l$, as in \eqref{ch11:cover-transfer}; it
must not be assigned the contraction factor for a vanishing extracted coefficient.
If its incoming mass is $\mathfrak m$, the other local terms with vanishing extracted coefficient have total mass $p_l$ with $2p_l\le\mathfrak m$
after increasing the lower bound on the reference coupling, and the large-field terms have normalized mass
$f_l\le\mathfrak m$. These are the termwise local bounds used
above, retaining the original events. Resolving real and imaginary
parts contributes at most a factor two, so
\[
 2(p_l+T_l\mathfrak m)+f_l
             \le(2L_l+2)\mathfrak m\le6L_l\mathfrak m
             \le L_l^4\mathfrak m.
\]
The last inequality holds for $L_l\ge3$. It applies in the
submultiplicative norm in the source variables to a whole incoming joint factor;
it does not multiply separate bounds for correlated insertions.
With $g$ groups left, the product of their remaining source weights is consequently
at least
\begin{equation}
 (a_j/a_l)^{4n}(\alpha_j/8)^{n(n-g)}\prod_iw_i.
 \label{ch18:group-weight}
\end{equation}
Each step loses at most the factor $L_l^{-4n}$, and a merger uses at most
$n$ of the weight-reduction factors for support enlargement. This proves the formula by
induction, including all data on supports and prescribed large-field sets.

\subsection{Extraction with a freely chosen residual weight}

Put
\[
 \delta_l^{\rm mark}=\widehat\beta_{l+1}^{-2},\quad
 \phi_l=\widehat\beta_{l+1}^{-1},\quad
 \Xi_l=e^{\sqrt{\widehat\beta_{l+1}}},\quad
 y_{t,l}=\widehat\beta_l^{-\nu_t}.
\]
Choose the fixed integer $\nu_t$ to absorb the finite polynomial and
attachment losses of the insertions at degree $t$. The stage lower bounds on the reference coupling ensure
$y_{t,l}C_l\le1$ and
$\phi_l L_l^{t+1}y_{t,l}\le y_{t,l+1}$.
These are finite requests on the lower bound on the reference coupling after the ratio is fixed.
All their logarithms, as well as $\log\Xi_l$, are $o(\tau_l)$.

For a fixed exterior reflection cell at scale $k$, let
\[
 k_* =\min\{l\ge k:B_l^2\le a_k/a_l\},\qquad
 N_k=\exp\{(a_k/a_{k_*})^4q_{k_*}^{\rm site}+5\}.
\]
The stopping-scale argument of
\eqref{ch13:sharp-schedule} gives $\tau_{k_*}/\tau_k\to1$
and $\log N_k=o(\tau_{k_*})$. Indeed if $k_*=l+1>k$,
minimality gives $a_k/a_{k_*}\le L_lB_l^2$; the last condition
of that equation bounds both the normalization exponent and
$\tau_{k_*}-\tau_k$ by $o(\tau_{k_*})$.
When $\tau_k=o(T)$ the same proof gives $\log N_k=o(T)$.

Suppose the cluster has become one group at scale $m+1$, with
product $\omega$ of its remaining source weights. Set
\[
 q_m^{\rm ext}=4\Xi_ma_{m+1}^t/\omega,\qquad
                    E_m=8(1+C_m)q_m^{\rm ext}.
\]
Taylor extraction through radial degree $t$ gives coefficient and
complement each bounded by $q_m^{\rm ext}$. To verify this, divide
by this product of weights, apply the integral Taylor formula in the
local derivative norm, and sum the finite coefficient conversions
and the polynomial factors in the support size with their positive margin in the exponent. Their
bounds are fixed stage constants and powers of logarithmic radii
and inverse thresholds, hence eventually at most $\Xi_m$.
The factor $a_{m+1}^t$ converts the dimensionful density with insertions to the
normalized degree-$t$ density. Large-field terms keep their
original normalized majorants. Extraction of the coefficient of
all insertions in a group is a linear operation on the exact density;
products of partial modes are retained before taking it. Each
term is assigned the union of its supports and original prescribed large-field set,
so this operation never divides by the exact-cover gas of the density without insertions.

Choose a residual weight $\varpi>0$ at $k_*$ and thereafter
$\varpi_{l+1}=\phi_l\varpi_l$. The two sufficient inequalities are
\begin{align}
 \varpi(a_m/a_{k_*})^{t+1}y_{t,m}/y_{t,k_*}
        &\le\delta_m^{\rm mark}\omega/(16\Xi_m),\notag\\
 32(1+C_m)\varpi a_{k_*}^{-t}E_m(a_m/a_{k_*})^{1/8}
        &\le y_{t,k_*}/2.
 \label{ch18:residual-weights}
\end{align}
The first absorbs the local polynomial coefficient bound and the second the bound for its signed transfer.
Indeed the recursion for these bounds gives
$\varpi_l\le\varpi(a_l/a_{k_*})^{t+1}y_{t,l}/y_{t,k_*}$.
Combining this with \eqref{ch18:signed-transport} and the dimensionful
factor $a_l^{-t}$, the ratio of required to allowed weight at
scale $l$ is at most its value at $k_*$ times
$(a_l/a_{k_*})^{7/8}\le1$. Thus no intermediate scale has been
left unchecked.

The extracted matched field has coefficient norm at scale $m$
at most $2E_m$. The interior residual has chessboard seminorm
at most $2N_k/\varpi$. To see the latter on arbitrarily many
copies, put one residual in every reflected cell before integrating.
Taking the coefficient of the required source monomial in the sum with insertions contributes the factor $\varpi^{-1}$ per
cell; the normalized integral bound \eqref{ch11:root-bound}
contributes the factor $N_k$ per cell. Apply this separately to the original local
difference and its compensating column, and use the chessboard
triangle inequality. The rooted component estimates and bounds
used in this calculation depend on the number of insertions per cell, not on the
number of cells. All the terms involving at least two insertions generated during transfer
are included in these sums. The estimate is therefore stronger
than a bound on one value of a correlation functional.

\begin{proposition}[Lattice cluster replacement]
\label{ch18:lattice-replacement}
Let a cluster contain $n$ fields of total degree $W$. Choose
$0\le\theta_-<\theta<\theta_+$ and scales $k,m,j$ with
$\tau_k/T\to\theta_-$, $\tau_m/T\to\theta$, and $\tau_j/T\to\theta_+$. The cluster lies inside the centered block of scale $m+1$,
its points lie on the scale-$j$ grid with separation at least
$B_j^2a_j$, and its dependence sets lie inside the scale-$k$
cell. Put $Z=W\theta_++4n(\theta_+-\theta)$.
For each fixed auxiliary degree $t$ there are output labels
$v_{N,p}^{\rm out}\in V_t$, on the actual histories, satisfying
\begin{align}
 |x_m(v_{N,p}^{\rm out})|_1
       &\preccurlyeq e^{(Z-t\theta)T},\notag\\
 \left\|\prod_iO_{N,p}(v_i,x_i)
                  -O_{N,p}(t,v_{N,p}^{\rm out};0)\right\|_{\mathcal C_k}
       &\preccurlyeq e^{[Z-(t+1/16)(\theta-\theta_-)]T}.
 \label{ch18:lattice-estimate}
\end{align}
The difference is interior. Constants and thresholds are uniform in
cutoff and period on the ordered tails and in the number of copies
of this fixed local cluster.
\end{proposition}
\begin{proof}
The weights of normalized local insertions in
\eqref{ch13:weight-size} give $\prod_iw_i=e^{-W\theta_+T+o(T)}$.
Keep the coefficients of all incoming source monomials above $j$; separation permits their
bounds as terms with insertions but does not assert statistical independence. The grouping
estimate gives $\omega=e^{-ZT+o(T)}$ at $m+1$, and therefore
$E_m=e^{(Z-t\theta)T+o(T)}$.
Set
\[
 \varpi=\exp\{[(t+1/16)(\theta-\theta_-)-Z]T\}.
\]
In the first inequality of \eqref{ch18:residual-weights}, the
left exponent is $-Z-15(\theta-\theta_-)/16$, strictly smaller
than the right exponent $-Z$. In the second, the full left
exponent is $-(\theta-\theta_-)/16<0$ while the right side
has exponent zero. Both strict margins absorb every subpower
factor. The extraction estimate and $\log N_k=o(T)$ give the
claim. All support shifts take place within the specified cell;
reflections and conjugations retain the same bounds.
\end{proof}

\subsection{A consequence for undecomposed correlations without phases}
\label{ch18:mixed-replacement}

Fix finite correlation reconstruction maps onto the filtered spaces $F_D$.
More explicitly, choose a basis in each successive quotient and
extend it to $F_D$, with dual coordinates realized by finite
linear combinations of separated probe correlations. Arrange that
a degree-$i$ coordinate vanishes on $F_{i-1}$, and that reconstruction
fixes every field in $F_D$. Keep each old coordinate functional as
its original correlation formula on every larger field space. When
lifting a new quotient basis, subtract its values under all the old
coordinates, so those coordinates annihilate the new basis vectors.
This is possible by finite-dimensional separation by all spectator
correlations and successive subtraction of lower degrees. Consequently
the degree-$D$ reconstruction of a field in $F_t$, $t>D$, is exactly its
low-degree component in this fixed biorthogonal system. Denote the resulting reconstruction of a word by
$\mathcal P_Dw$. It depends on the cluster, not on its later spectator.
No rotation average or phase measure is used here.

\begin{proposition}[Local replacement for undecomposed correlations]
\label{ch18:mixed-bound}
Let $x_i=o+r z_i$, with the $z_i$ in a fixed compact separated
configuration set, and let $W$ be the total cluster degree.
For any fixed separated exterior word $v$, any $D$ and $\varepsilon>0$,
\begin{align}
 |b_\alpha(w;o)|&\le C r^{|\alpha|-W-\varepsilon},\notag\\
 |S(wv)-S((\mathcal P_D(\tau_{-o}w))(o)v)|
                   &\le C r^{D+1-W-\varepsilon}.
 \label{ch18:mixed-product-rate}
\end{align}
Here $b_\alpha$, $|\alpha|\le D$, are the fixed probe coordinates. The same field
$\mathcal P_D(\tau_{-o}w)$ works for every exterior word. The
estimates are uniform on fixed bounded label sets and separated
geometries, and hold with any fixed finite derivatives in the positions,
with their increased scaling dimensions.
\end{proposition}
\begin{proof}
First use generic points of the scale-$j$ grid. Put $T=\log(1/r)$ and
choose $0<\theta_-<\theta<1<\theta_+$. The cluster then fits
the block and separation conditions for large $T$.
For each fixed geometry, the output labels in
\eqref{ch18:lattice-estimate} are bounded at fixed $m$ and hence
at the reference scale by the inverse-column bound. Take a
subsequence along which the coefficients converge, and call its field $A_t$.
The chessboard inequality bounds its residual paired with any
fixed exterior word by exponent
$Z-(t+1/16)(\theta-\theta_-)$. The same subsequence works for all
spectators, since the coefficient limit was taken before pairing.

The graded quotient estimate \eqref{ch18:graded-distance} shows
that a degree-$i$ coordinate of $A_t$ has exponent $Z-i\theta$.
Its coordinate functional vanishes on $F_{i-1}$; this is where
an ungraded coefficient estimate would lose the desired power.
Choose $t$ so large that
$(t+1/16)(\theta-\theta_-)\ge(D+1)\theta$.
This also bounds every raw coordinate being evaluated, whose degree is
at most $D$. We do not require the residual to bound all coordinates
through the auxiliary degree $t$. Pairing the residual
with the fixed probes bounds the coordinate errors. Reconstruct
$A_t$ and omit degrees $i>D$. Each omitted coordinate has exponent
at most $Z-(D+1)\theta$, and its fixed exterior correlations are
bounded. This gives that exponent for the raw replacement error,
and $Z-i\theta$ for each raw coordinate.

For completeness the grid restriction can be removed without
assuming differentiability of the selected coefficients of $A_t$.
Use the fixed correlation reconstruction instead. Its coefficients and
the pairings of its residual are smooth functions of the cluster positions.
Their derivatives replace an input by a field of degree one higher.
The sharp separated bounds, using fixed cells for exterior probes
and shrinking cells only for cluster factors, give a crude exponent
$W+1$: choose total loss one in the exponent. Thus one can take $c_0=1$.
The grid exponent is $\theta_+W+c_1$, where
$c_1=4n(\theta_+-\theta)-(D+1)\theta$ for the residual, and
$c_1=4n(\theta_+-\theta)-i\theta$ for a degree-$i$ coordinate.
These numbers are independent of $W$ with $n,D,i$ fixed.
Multiplicative constants may depend on each of the finitely many
higher derivative orders used in the following induction. For all sufficiently large total
degrees the crude estimate implies the grid target, because
$\theta_+>1$. Use downward induction on the total degree. Moving to a generic grid
configuration contributes the factor $O(a_j)=e^{-\theta_+T+o(T)}$; the derivative
bound at degree $W+1$ gives precisely the target exponent at
$W$. The connecting paths remain separated. Only finitely many
higher derivatives are needed for each starting degree, so their
constants and thresholds can be chosen together. Finally let
$\theta\uparrow1$ and $\theta_+\downarrow1$ and absorb the
remaining loss into the prescribed $\varepsilon$.
This proves the estimate for undecomposed correlations using neither phase symmetry nor
an operator-product theorem.
\end{proof}

